%! Setting Up --- Learning from Data — Unit 8, Lesson 8.0 — Guided Notes (Answer Key)
\documentclass[10pt]{article}
\usepackage{linalg-article}
\usepackage{linalg-key}   % pulls in -boxes; do NOT also load -boxes
\newcommand{\TallMath}[1]{$\displaystyle #1\rule[-1.4em]{0pt}{3.2em}$}
\newcommand{\vv}{\mathbf{v}}
\newcommand{\bb}{\mathbf{b}}
\newcommand{\zero}{\mathbf{0}}
\newcommand{\relu}{\mathrm{ReLU}}

\begin{document}

\pageheader{Unit 8, Lesson 8.0}{Guided Notes --- Answer Key}
\namedateperiod

\vspace{0.12in}

\begin{objectivebox}
\textbf{By the end of this lesson, I will be able to\dots}
\begin{itemize}\setlength{\itemsep}{2pt}
  \item explain that a \emph{learning function} (a neural network) is built by alternating a \emph{linear step} (multiply by a weight matrix, add a bias) with a \emph{nonlinear step} (the ReLU);
  \item evaluate $\relu(x)=\max(0,x)$ and compute one layer $\relu(A\vv+\bb)$;
  \item explain how stacking ReLUs builds \emph{piecewise-linear} functions that bend to fit data, and that \emph{learning} means adjusting the weights to make a \emph{loss} smaller.
\end{itemize}
\end{objectivebox}

\vspace{0.08in}

\begin{vocabbox}
\small
Fill in each term as we build it together.
\vspace{2pt}
\vocabans{Learning function $F(\vv)$ (a neural network)}{a function that turns an input vector into a prediction, built by alternating linear steps ($A\vv+\bb$) with the nonlinear ReLU}
\vocabans{ReLU: $\relu(x)=\max(0,x)$ (the nonlinear step)}{keeps a positive number, replaces a negative with $0$; its single fold is what lets a network bend}
\vocabans{Weights: a weight matrix $A$ and a bias $\bb$}{the numbers a network gets to choose; the entries of $A$ scale/mix the inputs and $\bb$ shifts them}
\vocabans{Piecewise-linear function (straight pieces joined at folds)}{a function made of straight segments meeting at corners; sums of ReLUs produce them, and they can trace any bending data}
\vocabans{Loss, and gradient descent}{loss $=$ sum of squared errors between predictions and targets; gradient descent nudges the weights ``downhill'' to make the loss smaller}
\end{vocabbox}

\begin{hookbox}
For seven units we used matrices to \emph{move}, \emph{project}, \emph{stretch}, and \emph{factor}
the plane. The surprise of Unit~8: those same matrix tools are the engine inside modern
\textbf{artificial intelligence}. A neural network is just a \emph{function} that turns an input into
a prediction --- and it is built almost entirely from things you already know.
\begin{itemize}\setlength{\itemsep}{1pt}
  \item In Unit~4 you found the \emph{best-fit line} through data. What were the two numbers you got to choose? \ansline{The slope $w$ and the intercept $b$ of the line $y=wx+b$ --- its two ``weights.''}
  \item A straight line can only go straight. If the data \emph{bends}, what would a good fitting function need to be able to do? \ansline{It needs to \emph{bend} too --- change slope, not stay a single straight line.}
\end{itemize}
\end{hookbox}

\begin{notesbox}{1. Learning from Data = Finding a Function That Fits}
\textbf{Learning from data} means: given input--output examples, find a \emph{function} that reproduces
them and predicts new ones. The simplest learning function is the Unit~4 \textbf{best-fit line}
$y=wx+b$ --- it has just two \emph{weights} to choose: the slope $w$ and the intercept $b$. Fitting the
data means picking $w,b$ to pass as close as possible to the points.

\vspace{2pt}
But a line can only go \emph{straight}. Real data bends, so we need functions that can \textbf{bend} too
--- while still being built from simple straight pieces. That one new ability is the whole idea of
today's lesson: the ellipse of Unit~7 gave way to \emph{learning}, and the tool that lets a function
bend is a single, very simple function.
\end{notesbox}

\begin{notesbox}{2. The One New Ingredient: ReLU $=\max(0,x)$}
Here is the only new function in Unit~8. The \textbf{ReLU} (``rectified linear unit'') keeps a positive
number and replaces a negative number with $0$:
\[
  \relu(x)=\max(0,x).
\]
Fill in:
\[
  \relu(3)={\color{keyred}\mathbf{3}},\qquad \relu(-2)={\color{keyred}\mathbf{0}},\qquad \relu(0)={\color{keyred}\mathbf{0}},\qquad \relu(5)={\color{keyred}\mathbf{5}}.
\]
As a graph it is flat at $0$ for negative inputs, then rises along the line $y=x$: a single \textbf{fold}
at the origin. That one bend is the secret. Add and shift a few ReLUs and their folds combine into a
\textbf{piecewise-linear} function --- straight pieces joined at corners --- that can bend to trace
\emph{any} data. For example, $F(x)=\relu(x)-\relu(x-2)$:
\[
  F(-1)={\color{keyred}\mathbf{0}},\qquad F(1)={\color{keyred}\mathbf{1}},\qquad F(3)={\color{keyred}\mathbf{2}}.
\]
This $F$ rises from $0$ to $2$ over $0\le x\le 2$, then stays flat --- two folds, built from two ReLUs.
That is Lesson~8.1.

\begin{center}
\begin{tikzpicture}[scale=0.62]
  % LEFT: one ReLU = one fold
  \begin{scope}[shift={(-3.0,0)}]
    \draw[->, charcoal, thin] (-2.2,0)--(2.5,0) node[right]{\scriptsize $x$};
    \draw[->, charcoal, thin] (0,-0.5)--(0,2.6) node[above]{\scriptsize $\relu(x)$};
    \draw[ultra thick, burgundy] (-2,0)--(0,0)--(2,2);
    \fill[burgundy] (0,0) circle (1.6pt);
    \node[charcoal] at (0.9,-0.9) {\scriptsize one fold};
  \end{scope}
  % ARROW
  \draw[->, very thick, goldacc] (-0.2,1.0)--(1.7,1.0) node[midway, above]{\scriptsize combine};
  % RIGHT: several ReLUs = a bent fit through data
  \begin{scope}[shift={(4.0,0)}]
    \draw[->, charcoal, thin] (-0.4,0)--(4.4,0) node[right]{\scriptsize $x$};
    \draw[->, charcoal, thin] (0,-0.5)--(0,2.9) node[above]{\scriptsize $y$};
    \fill[charcoal] (0.5,0.4) circle (1.6pt);
    \fill[charcoal] (1.3,1.3) circle (1.6pt);
    \fill[charcoal] (2.1,2.1) circle (1.6pt);
    \fill[charcoal] (2.9,1.7) circle (1.6pt);
    \fill[charcoal] (3.7,1.1) circle (1.6pt);
    \draw[ultra thick, royalblue] (0.3,0.2)--(2.0,2.0)--(3.8,1.0);
    \node[charcoal] at (2.0,-0.9) {\scriptsize a few ReLUs $\to$ a bent fit};
  \end{scope}
\end{tikzpicture}
\end{center}
\end{notesbox}

\begin{notesbox}{3. A Learning Function Is a Chain of Matrix Steps}
A neural network takes an input vector $\vv$ and pushes it through \textbf{layers}. Each layer does two
things: a \emph{linear step} --- multiply by a \textbf{weight matrix} $A$ and add a \textbf{bias} vector
$\bb$ (Unit~1) --- then the \emph{nonlinear step} --- apply ReLU to each entry:
\[
  \text{one layer:}\qquad \vv \;\longmapsto\; \relu(A\vv+\bb).
\]
Take $A=\begin{bsmallmatrix} 1 & -1\\ 1 & 1 \end{bsmallmatrix}$, $\bb=\begin{bsmallmatrix} 0\\ 0 \end{bsmallmatrix}$,
and the input $\vv=\begin{bsmallmatrix} 1\\ 2 \end{bsmallmatrix}$. First the linear step:
\[
  A\vv+\bb=\begin{bmatrix} 1 & -1\\ 1 & 1 \end{bmatrix}\begin{bmatrix} 1\\ 2 \end{bmatrix}
  =\begin{bmatrix}{\color{keyred}\mathbf{-1}}\\[2pt]{\color{keyred}\mathbf{3}}\end{bmatrix}.
\]
Then ReLU on each entry (keep positives, zero out negatives):
\[
  \relu\!\begin{bmatrix}{\color{keyred}\mathbf{-1}}\\[2pt]{\color{keyred}\mathbf{3}}\end{bmatrix}
  =\begin{bmatrix}{\color{keyred}\mathbf{0}}\\[2pt]{\color{keyred}\mathbf{3}}\end{bmatrix}.
\]
The ReLU turned the negative entry into \ans{$0$}. \textbf{Matrix multiplication does all the heavy
lifting} in every layer --- a neural network is mostly the Unit~1 operation $A\vv$, repeated, with a ReLU
between the steps. The bias $\bb$ just \emph{shifts} where each fold sits.
\end{notesbox}

\begin{notesbox}{4. How It ``Learns'' --- and Why It Matters}
The weights (the entries of every $A$ and $\bb$) start as guesses. To \textbf{learn}, we measure how wrong
the predictions are with a \textbf{loss} --- the Unit~4 \emph{sum of squared errors} between the network's
outputs and the true targets. If the network predicts $3,5,4$ and the targets are $2,5,6$:
\[
  \text{loss}=(3-2)^2+(5-5)^2+(4-6)^2={\color{keyred}\mathbf{1}}+{\color{keyred}\mathbf{0}}+{\color{keyred}\mathbf{4}}={\color{keyred}\mathbf{5}}.
\]
Learning means nudging the weights to make this loss \emph{smaller} --- rolling downhill on the loss until
it bottoms out. That downhill method is \textbf{gradient descent} (Lesson~8.3). The shape of the data
itself --- its \emph{mean}, \emph{variance}, and \emph{covariance} --- is described by the statistics of
Lesson~8.4.

\vspace{2pt}
\textbf{Why it matters.} A neural network is not new mathematics --- it is \ans{matrix} multiplication
(Unit~1) plus one simple bend (ReLU), trained by least squares (Unit~4). The linear algebra you already
know is the engine of modern AI.

\vspace{2pt}
\textbf{The road ahead.} \textbf{8.1} piecewise-linear learning functions\; $\bullet$\;
\textbf{8.2} convolutional neural nets\; $\bullet$\; \textbf{8.3} minimizing loss by gradient descent\;
$\bullet$\; \textbf{8.4} mean, variance, and covariance.
\end{notesbox}

\begin{practicebox}
Show your work.
\begin{enumerate}[leftmargin=1.6em, itemsep=4pt]
  \item Evaluate $\relu(7)$, $\relu(-4)$, and $\relu(0)$. \ansline{$\relu(7)=7$, $\relu(-4)=0$, $\relu(0)=0$.}
  \item For $A=\begin{bsmallmatrix} 2 & -1\\ -1 & 2 \end{bsmallmatrix}$ and $\vv=\begin{bsmallmatrix} 1\\ 3 \end{bsmallmatrix}$, compute one layer $\relu(A\vv)$ (take $\bb=\zero$). \ansline{$A\vv=\begin{bsmallmatrix} -1\\ 5 \end{bsmallmatrix}$, so $\relu(A\vv)=\begin{bsmallmatrix} 0\\ 5 \end{bsmallmatrix}$ (the $-1$ becomes $0$).}
  \item A network predicts $4,2$ while the targets are $3,4$. Compute the loss (sum of squared errors). \ansline{$(4-3)^2+(2-4)^2=1+4=5$.}
\end{enumerate}
\end{practicebox}

\begin{teachernote}
The spine is Section~2 (ReLU $=\max(0,x)$, the one nonlinear ingredient) and Section~3 (one layer
$\relu(A\vv+\bb)$ --- matrix multiplication, Unit~1, does the work). Section~1 frames ``learning $=$
fitting a function'' off the Unit~4 best-fit line and motivates \emph{bending}; Section~4 names the loss
(Unit~4 sum of squares) and gradient descent, and delivers the punchline: a neural net is the linear
algebra they already know. Keep the arithmetic light --- this is an on-ramp, not a computation drill.
Most common slips: applying ReLU to the \emph{whole} vector instead of \emph{each entry}; thinking ReLU
changes positive numbers; forgetting to square before summing in the loss. Sentence to land: matrix
multiply, then bend with ReLU, then make the loss smaller --- that is a neural network learning.
\end{teachernote}

\end{document}
